\section{A direct geometric bound}
\label{v10:geometry}
The moment argument can be made without a succession of feedback updates.
The useful quantity is a single cost for changing the sign probability,
moving the inner support endpoints, and adding conditional noise. We prove
its scalar consequences directly, then evaluate it once on each of the
already available probability intervals.

Let $X$ have mean zero, variance one and compact support with an actual
central gap $(-A,B)$, where $A,B>0$ are support endpoints. Set
\[
 p=\mathbb P(X>0),\quad q=1-p,\quad v=pq,\quad
 m=\mathbb EX_+=\mathbb EX_-,\quad \bar A=m/q,\quad\bar B=m/p,
\]
\[
 \theta=1-m^2/v,\quad w=1-\theta,\quad
 \ell=m/v=\sqrt{w/v},\quad d=1-(A+B)/\ell.
\]
There is no atom at zero. Conditional means give $A\le\bar A=\ell p$,
$B\le\bar B=\ell q$, so $0\le d<1$. Write
$W=X-\mathbb E[X\mid\mathbf1_{X>0}]$ and define measures on magnitudes by
\[
 \int f\,d\mu_\pm=\mathbb E[W^2\mathbf1_{\{\pm X>0\}}f(|X|)].
\]
Their total mass is $\theta$. For $\sigma=\sqrt{10}$ put
\[
 Q(p)=(\sigma+1)p^2+(\sigma-1)q^2-6pq
      =(6+2\sigma)(p-p_*)^2,\qquad p_*=(4-\sigma)/2.
\]
Let $\rho=\mathbb E|X|^3$, $\gamma=\mathbb EX^3$,
$\beta=\mathbb E[X|X|]$,
$J=\sum_{(a,b)}\tau_{a,b}(b-a)^2$, and
$D=(3+\sigma)\rho-\gamma-3J$. The sum is over bounded support gaps and
$\tau(t)=\mathbb E[X\mathbf1_{X\ge t}]$ is constant on each gap.

\begin{proposition}[One common geometric cost]\label{v10:geometry:cost}
With the preceding definitions,
\begin{equation}\label{v10:geometry:direct}
 D\ge\ell\{wQ(p)+6wd(1-d)+\theta K(p,d)\},
\end{equation}
where
\[
 K(p,d)=\min\{3(\sigma+1)q-(\sigma+5)d,
                    3(\sigma+3)p-(\sigma+7)d\}.
\]
No maximizing-contact or predecessor hypothesis is needed for this
geometric inequality.
\end{proposition}
\begin{proof}
For the positive outer gaps,
\[
 J_o^+\le2\int_B^\infty(t-B)\tau(t)\,dt
       =\mathbb E[X_+(X_+-B)^2],
\]
and the same inequality holds on the negative side. Endpoint atoms are
included; isolated integration endpoints have zero Lebesgue measure.
Add the central gap contribution $m(A+B)^2$. Substitution gives
\[
 D\ge \ell wQ(p)+E_0+\int k_+\,d\mu_++\int k_-\,d\mu_-,
\]
\[
 E_0=6m\{(A+B)(\bar A-A+\bar B-B)+(\bar A-A)(\bar B-B)\},
\]
\[
 k_+(x)=(\sigma-1)(x+2\bar B)+6B,\qquad
 k_-(x)=(\sigma+1)(x+2\bar A)+6A.
\]
For completeness, the substitutions use
$\mathbb EX_+^3=m\bar B^2+\int(x+2\bar B)d\mu_+$ and its negative
analogue, together with
$\mathbb EX_+^2-mB=\mu_+(\mathbb R)+m(\bar B-B)$.
Put $a_0=(\bar A-A)/\ell$, $a_1=(\bar B-B)/\ell$. Then
$a_0,a_1\ge0$, $a_0+a_1=d$, and
\[
 E_0/\ell=6w\{d(1-d)+a_0a_1\}.
\]
On the positive measure $x\ge B$, so
$k_+(x)/\ell\ge3(\sigma+1)q-(\sigma+5)a_1$.
On the negative measure the corresponding lower bound is
$3(\sigma+3)p-(\sigma+7)a_0$.
Drop $a_0a_1$, use $a_i\le d$, and sum the two nonnegative measures.
Their combined mass is $\theta$, proving the result.
\end{proof}

\subsection{Direct global localization}
The following numerical premises are the coarse entry's global output.
They are not conclusions of the older moment iteration:
\begin{equation}\label{v10:geometry:input}
 p\in[.24,.54],\quad n\ge42,\quad nD\le32.20665378,
 \qquad |X|\le157/50.
\end{equation}
Only the final third-moment bound uses the support radius. Define
\[
 D_0=.76682509,\quad s=3.162277660,\quad s^+=3.162277661,
 \quad z=(s-1-D_0)/3,
\]
\[
 T=7/125=.056,\quad W_0=1-T=.944,\quad d_0=17/250=.068.
\]
All displayed decimals in this section denote exact rationals.

\begin{theorem}\label{v10:geometry:global}
Under \eqref{v10:geometry:input},
\[
 p>1/4,\qquad \theta<.056,\qquad d<.068,\qquad
 \lambda:=n\theta/\ell^2<.616,\qquad \rho<1.426.
\]
The numerical part reduces to four quartics, two cubics and elementary
rational inequalities specified below.
\end{theorem}
\begin{proof}
The preceding outer-gap bound also gives $J\le\rho+2mAB$. Thus
\[
 D\ge\sigma\rho-\gamma-6mAB\ge\sigma-1-3AB,
\]
using $\rho\ge1$, $\gamma\le\rho$, $m\le1/2$. Consequently
\begin{equation}\label{v10:geometry:domain}
 AB\ge z>0,\quad w\ge AB\ge z,\qquad
 0\le d\le1-2\sqrt{zv/w}=:d_{\max}.
\end{equation}
Here $AB\le\bar A\bar B=w$ and $A+B\ge2\sqrt{AB}$.

First, $\theta<1/5$ follows without iteration. Since
$B\ge z/\bar A$, the positive coefficient in the common budget satisfies
\[
 k_+/\ell\ge q\{(\sigma+5)z/w+2(\sigma-1)\}\ge3.13.
\]
The negative coefficient is bounded by
$p\{(\sigma+7)z/w+2(\sigma+1)\}\ge3.13$ in the same direction.
These last bounds use $w\le1$, $q\ge.46$, $p\ge.24$.
Therefore $D\ge3.13\ell\theta$; since $\ell\ge2\sqrt z$, a value
$\theta\ge1/5$ would give $D\ge(313/250)\sqrt z>D_0$.

Suppose next $T\le\theta<1/5$. For fixed $p,\theta$, each branch
of \eqref{v10:geometry:direct} is concave in $d$. It suffices to
exclude $d=0$ and $d=d_{\max}$. At the latter endpoint put
$g=1-d_{\max}=2\sqrt{zv/w}$. Exact rational square comparisons give
$.59<g<.77$, $g(1-g)>.17$ and $\ell>1.7$.
Both noise coefficients are positive, since
$3(\sigma+1)q,3(\sigma+3)p>4.4$ and
$4.4-(s^++7)\cdot.41>0$. Hence the inward cost alone exceeds
$1.7\cdot6\cdot.8\cdot.17>D_0$.
At $d=0$ the two branches are
\[
 F_C(\theta)=\sqrt{(1-\theta)/v}\{(1-\theta)Q(p)+\theta C(p)\},
\]
where $C(p)=3(\sigma+1)q$ or $3(\sigma+3)p$.
Their derivative has the sign of
$(2-3\theta)C-3(1-\theta)Q$, which is positive because
$Q\le.4$, $C>4.4$, $\theta\le.2$.
It remains to prove $F_C(T)>D_0$.

Here and below replace $\sigma$ by its rational lower bound $s$ in
$Q,C$, writing $Q_s,C_s$. The required positive-base squared inequalities
are
\[
 P_\pm(p)=W_0\{W_0Q_s(p)+TC_{s,\pm}(p)\}^2-D_0^2p(1-p)>0.
\]
They have $P_\pm''\ge2D_0^2$ on $[.24,.54]$. If a twice differentiable
function satisfies $P''\ge2a>0$, then
\begin{equation}\label{v10:geometry:tangent}
 P(p)\ge P(x)-P'(x)^2/(4a).
\end{equation}
This follows by integrating the second derivative and completing the square.
The exact evaluations at $x=.45,.39$ give respectively lower bounds
$.0017$ and $.011$. This excludes both endpoints and proves $\theta<T$.

From \eqref{v10:geometry:domain}, $d<.42$. Each branch of the direct
bound now increases in $d$, since
$6W_0(1-2\cdot.42)-T(s^++7)>0$.
At $d=d_0$, it also increases in $\theta$: the derivative numerator is
\[
 (2-3\theta)(C-(\sigma+5\text{ or }\sigma+7)d_0)
 -3(1-\theta)\{Q+6d_0(1-d_0)\}>0.
\]
The least value is therefore at $\theta=0$. The remaining inequality is
\[
 P_d(p)=\{Q_s(p)+6d_0(1-d_0)\}^2-D_0^2p(1-p)>0.
\]
Again $P_d''\ge2D_0^2$; formula~\eqref{v10:geometry:tangent} at $.424$
gives a lower bound $.0012$. Thus $d<d_0$.

Since the direct bound increases in $d$, it gives
$D\ge\ell\{wQ+\theta C\}$ for at least one of the two values of $C$.
Multiplying by $n$ and using $n\theta=w\lambda/v$ yields
\[
 \lambda\le\max\left\{
 \frac{p[17.558-42Q(p)]}{3(\sigma+1)},
 \frac{q[17.558-42Q(p)]}{3(\sigma+3)}\right\}.
\]
Here $\ell w\ge2W_0^{3/2}$ and
$32.20665378/(2W_0^{3/2})<17.558$.
The numerators are positive since $Q\le.4$. Replacing $Q,\sigma$ by
their lower enclosures gives rational cubics
$U_+,U_-$. The positive cubic increases up to $p=1/3$ and satisfies
$U_+''\le-13$ thereafter. The negative cubic has $U_-''\le-12$
on the entire interval. The concave counterpart of
\eqref{v10:geometry:tangent}, evaluated at $.455$ and $.392$, gives
$U_+<.615289$, $U_-<.565358$, proving the stated bound on $\lambda$.

For the third moment, $\ell>1.94$, $B>1.94(.46-.068)$ and
$\bar B<1.8$. Thus, for $x\le3.14$,
\[
 \frac{k_+(x)}{x+2\bar B}
 \ge s-1+\frac{6\cdot1.94(.46-.068)}{3.14+3.6}>5/2;
\]
the negative ratio is already at least $\sigma+1>5/2$.
The exact third-moment identity and common budget imply
\[
 \rho\le\frac{w^{3/2}}{\sqrt v}\{1-2v-\tfrac25Q(p)\}+\tfrac25D_0
 \le\frac{1-2v-\tfrac25Q_s(p)}{\sqrt v}+\tfrac25D_0.
\]
The bracket is positive since $Q\le.4$. For
$M=1.426-\tfrac25D_0$, the final quartic is
\[
 P_\rho(p)=M^2p(1-p)-\{1-2p(1-p)-\tfrac25Q_s(p)\}^2.
\]
It satisfies $P_\rho''\ge.4$ on the full interval, and
\eqref{v10:geometry:tangent} at $.258$, with $a=.2$, gives
$P_\rho>.00015$. Positivity of both squared bases proves $\rho<1.426$.
Finally $D\ge Q(p)/\sqrt{pq}$ follows by the already proved monotonicities
in $d$ and $\theta$. This ratio decreases on $[.24,.25]$, and its value
at $.25$ exceeds $D_0$ by an exact squared comparison. Hence $p>1/4$.
\end{proof}

\paragraph{Finite arithmetic.}
The four quartics are specified by their displayed formulas, not by fitted
coefficients. Their second derivatives are quadratic, whose minima on a
closed interval occur at endpoints or at the eligible vertex. Cubic second
derivatives are affine. Thus the stated derivative bounds and tangent
values require only rational evaluations. The accompanying pure-Python
checker reconstructs the polynomials and performs these operations with
\texttt{Fraction}. It also checks positivity before squaring and the
rational enclosures $s^2<10<(s^+)^2$. These checks prove precisely the
scalar inequalities used above, conditional on the analytic derivation.

\subsection{One evaluation on each probability interval}
The global lemma also removes the repeated updates needed to obtain the
existing support strips. Let a closed entry panel be $I=[l,u]$, with
$D\le D_I$ and $nD\le42D_I$. Put $v_-,v_+$ for its actual minimum and
maximum of $p(1-p)$, including $1/4$ when the panel crosses $1/2$.
Let $Q_I=\max(0,\min_I Q_s)$ and
\[
 C_I=\min\{3(s+1)(1-u),3(s+3)l\},\qquad e_0=\sqrt{W_0/v_+}.
\]
Since $D\ge\ell\{Q+\theta(C-Q)\}$ and $\ell\ge e_0$, a negative
$D_I/e_0-Q_I$ excludes the whole panel. Otherwise one evaluation gives
\begin{equation}\label{v10:geometry:panel}
 \Theta_I=\min\left(T,\frac{D_I/e_0-Q_I}{C_I-Q_I}\right),\quad
 e_I=\sqrt{(1-\Theta_I)/v_+},\quad
 J_I=\frac{D_I}{e_I(1-\Theta_I)}-Q_I,
\end{equation}
\[
 d_I=\min\left(d_0,\frac{1-\sqrt{1-2J_I/3}}2\right),\quad
 A_I=e_I(l-d_I),\quad B_I=e_I(1-u-d_I).
\]
All roots and divisions are checked on their stated domains. If $J_I<0$
the panel is impossible; otherwise the smaller quadratic root is correct
because the global result has already proved $d<.068<1/2$.
These inequalities follow from
$D\ge\ell w\{Q+6d(1-d)\}$, since the noise coefficients remain positive.

For a coherent third-moment estimate define
\[
 \bar A_u=\sqrt{u/(1-u)},\quad\bar B_u=\sqrt{(1-l)/l},\quad
 L_I=R+2\max(\bar A_u,\bar B_u),
\]
\[
 k_I=\min\{s-1+6B_I/(R+2\bar B_u),\;
                 s+1+6A_I/(R+2\bar A_u)\},\quad
 r_I=D_I-e_I(1-\Theta_I)Q_I.
\]
The noise term in $\rho$ is at most $\min(L_I\theta,r_I/k_I)$.
Write $r_B=(1-2v_-)/\sqrt{v_-}$ and
$\alpha_I=(3/2)\sqrt{1-\Theta_I}$. Integrating the derivative of
$(1-\theta)^{3/2}$ gives the affine bound
$1-\alpha_I\theta$, which stays positive since
$\alpha_I\theta\le(3/2)T<1$. Hence
\[
 \rho\le\min\left(1.426,\;
 r_B+(L_I-\alpha_Ir_B)_+
             \min\{\Theta_I,r_I/(k_IL_I)\}\right).
\]
This is the maximum of two affine pieces meeting at their turnover; it
does not combine incompatible allocations of the same noise.
The other exports are
\[
 1-2u-2(1-u)\Theta_I\le\beta\le1-2l+2l\Theta_I,
 \qquad e_I\le\ell\le1/\sqrt{v_-},
\]
\[
 \lambda\le\frac{v_+}{C_I}
      \left\{\frac{42D_I}{e_I(1-\Theta_I)}-42Q_I\right\}.
\]
The same calculation retains a joint displacement--noise constraint. Set
\[
 a_I^*=6(1-d_I)-\frac{\Theta_I}{1-\Theta_I}(s^++7)>0.
\]
In the direct budget,
$6wd(1-d)-\theta(\sigma+5\text{ or }\sigma+7)d\ge w a_I^*d$.
Multiply by $n$, use $n\theta=w\lambda/v$, $n\ge42$ and
$\ell w\ge e_I(1-\Theta_I)$. All remaining costs are nonnegative, so
\begin{equation}\label{v10:geometry:joint}
 a_I^*d+\frac{C_I\lambda}{42v_+}\le J_I.
\end{equation}
If $J_I=0$ this forces $d=\lambda=0$; otherwise it is a single shared
budget with intercepts $J_I/a_I^*$ and $42v_+J_I/C_I$. It does not permit
spending the full reserve independently on displacement and noise.
A second useful line follows by bounding the same mixed cost in the other
direction, using $d\le d_I$. Averaging the two valid lower bounds gives
\[
 \widehat a_I d+\frac{\widehat C_I\lambda}{42v_+}\le J_I,\qquad
 \widehat a_I=6(1-d_I)-\frac{\Theta_I(s^++7)}{2(1-\Theta_I)},\quad
 \widehat C_I=C_I-\frac{(s^++7)d_I}{2}.
\]
Both coefficients are checked positive. This uses
$2(\theta/w)d\le[\Theta_I/(1-\Theta_I)]d+d_I(\theta/w)$, so it averages
valid bounds for the same cost rather than adding two expenditures.
All these are direct formulas; no result is fed back into an earlier formula.

\begin{proposition}[Connection with the current entry]
For the thirty closed nonempty panels of the selected v7 entry, the preceding
single evaluation excludes panels 24 and 53 and retains panels 25 through 52.
Under the fixed-target contact and sharp-predecessor hypotheses of the
retained heat/ray lemma, outward rational evaluation supplies all thirteen existing
$|U|<.6$ strips on $p\in[.40,.53]$ through the retained all-threshold
heat and ray inequalities. On $[.41,.42]$ it gives
\[
 \theta\le.049086874003,\quad d\le.066505471827,\quad
 \rho\le1.204660335428,\quad\lambda\le.526219192656,
\]
\[
 1.975749037978\le\ell\le2.033209198406,
 \qquad .103059226156\le\beta\le.220251236683.
\]
The averaged joint bound has outward intercepts $.069772925988$ and
$.551973904179$, both smaller than the old central moment resource's
$.0730144376$ and $.5728987925$. These imply every outward moment, noise
and scale premise used by the
current central-panel confinement theorem. The raw-threshold and signed
phase bounds still come from their original separate predecessors.
\end{proposition}
\begin{proof}
The direct formulas above prove the moment and scale bounds on each entire
closed panel. The finite check reads only the named entry's frozen input,
retains all thirty nonempty panels, and uses directed rational square roots.
For the strip test put $t_-=e_I(l+.6)$, $t_+=e_I(1-u+.6)$, with $r$ the
proved rho cap and $b=-\beta_l$ or $b=\beta_u$, respectively. The retained
all-threshold contact lemma bounds the contact expression at $t>\sqrt3$ by
\[
 \frac25 S^3\left\{\frac{t^3-3t}{6}
                +\frac{|t^4-6t^2+3|}{72}\right\}
 +S^3\left\{\frac{t^5}{4100}+\frac{3t}{820}
                 +\frac{17(3t^2-2)}{78720}\right\}+c_0Q_r(t),
\]
where $S=81/80$, $c_0=819/2000$, $e=8/875$ and
$Q_r(t)=-t^3+3rt^2/2+3bt+er$. At both boundaries the finite check requires
this expression negative and
\[
 Q_r,Q_r',Q_r''<0,\qquad
 L=\tfrac25(t^2/2+1/2+e)+c_0Q_r'<0,\quad L',L''<0.
\]
These are the existing whole-ray sufficient conditions, with their original
index, contact and predecessor hypotheses. Both sides pass exactly on panels
40 through 52. All thirteen closed intervals and both endpoints are retained;
the other fifteen panels are not claimed to satisfy these strip tests.
The displayed central vector consists of outward bounds from its single
rational evaluation; direct comparisons with the previous central theorem's
premises give the asserted connection.
\end{proof}

\paragraph{What this replaces.}
The eight-round moment iteration is unnecessary for the four global
bounds, the probability exclusions, the thirteen incoming strips and the
central-panel input vector just listed. The actual v7 entry, its support
radius and weighted $nD$ bound, the common gap budget and the heat/ray
lemma remain dependencies. This is not a proof of the remaining
high-noise panels, the count comparison or the finite base. Any downstream
calculation using individual old panel bounds requires a new input-by-input
check; the global improvements alone do not authorize such substitutions.
